\documentclass[11pt]{article}
\usepackage[margin=25mm]{geometry}
\usepackage{amsmath,amssymb,booktabs,hyperref}
\title{Eight bounded rank-four hyperbolic quotient searches}
\author{Independent recovery of fixed-source Table 3}
\date{}
\begin{document}\maketitle
\begin{abstract}
We certify finite marked quotients for eight diagrams appearing in Table 3
of Fernandes--Leemans--Weiss. Complete bounded normal-subgroup searches,
followed by intersection and actual coset-incidence tests, give 59 positive
marked kernels. We distinguish universal spherical residues from residues
whose diagrams are finite. In particular no regular hypertope with either
local convention realizes the printed order $336$ for $(3,4,3,4)$ or order $2160$
for $(3,5,3,4)$. Exact named extension descriptions and the ninth diagram
$(3,5,3,5)$ are outside this component.
\end{abstract}

\section{Conventions and the finite assertion}
Write $I=\{1,2,3,4\}$. For a symmetric Coxeter matrix $M$, put
\[
 W(M)=\langle s_i\ (i\in I)\mid s_i^2=1,
                     (s_i s_j)^{M_{ij}}=1\ (i<j)\rangle.
\]
In a square $(p,q,r,s)$ the edges $12,23,34,41$ have those labels,
and the diagonals have label2. In a string $[p,q,r]$ the edges are
$12,23,34$. For the fork $[5,3_3]$ the edges $12,13,14$ have
labels $5,3,3$. Unlisted off-diagonal entries are2.
These conventions agree, up to relabelling, with the native diagrams on
source physical page 12 and the types on physical page 14.

For a quotient $H=\langle t_1,t_2,t_3,t_4\rangle$, let
$H_J=\langle t_j:j\in J\rangle$ and $P_i=H_{I\setminus\{i\}}$.
The \emph{broad} column below requires a regular hypertope with this
exact diagram; its rank-three residues have finite diagrams. The
\emph{universal} column additionally requires each $P_i$ to be the full
finite Coxeter group for the corresponding rank-three subdiagram.
This distinction matters: the fixed source defines spherical by a finite
diagram on page 5, but explicitly restricts to spherical, rather than
projective, local geometries in its euclidean discussion on page 10.
We do not silently resolve that terminological tension for Section 6.

\begin{center}\small
\begin{tabular}{lrrrrr}\toprule
Diagram & Bound & Kernels & Positive & Broad minimum & Universal minimum\\\midrule
$(3,3,3,4)$&7680&14&10&168&1920\\
$(3,3,3,5)$&7200&5&2&2520&7200\\
$(3,4,3,4)$&9216&45&21&168&1152\\
$(3,5,3,4)$&2160&8&1&1080&$>2160$\\
$[3,5,3]$&8160&5&2&660&8160\\
$[4,3,5]$&7680&17&7&1920&3840\\
$[5,3,5]$&8160&6&3&3420&7200\\
$[5,3_3]$&7680&20&13&960&3840\\\bottomrule
\end{tabular}
\end{center}
Here kernels include the whole group (index1). Positive means distinct
normal kernels for the fixed ordered presentation, not isomorphism classes
after diagram automorphisms. There are120 kernels in total and59 positives.
The entry $>2160$ is only an exclusion through that bound, not an existence
or minimum assertion beyond it. All displayed attained minima have explicit
permutation witnesses. The original source's group names are not inferred
from orders or from descriptive strings printed by GAP.

\section{A complete coset test}
Suppose the involutions $t_i$ generate a finite group $H$, have exact pair
orders $M_{ij}$, and satisfy
\begin{equation}\label{ip}
 H_A\cap H_B=H_{A\cap B}\qquad(A,B\subseteq I).
\end{equation}
Form the incidence system whose objects of type $i$ are right cosets
$P_i g$, with incidence given by nonempty intersection. For
$J\subseteq I\setminus\{k\}$ let $E(J,k)$ be the set of type-$k$ objects
incident with every base object $P_i$, $i\in J$. The following exact test
is applied for all 32 choices:
\begin{equation}\label{flags}
 E(J,k)=\{P_k h:h\in H_{I\setminus J}\}.
\end{equation}
The GAP implementation compares the actual sets in the right-coset
permutation action. The independent integer-permutation implementation
constructs every right coset as a literal set and checks the equivalent
cardinality
$|E(J,k)|=|H_{I\setminus J}|/|H_{I\setminus(J\cup\{k\})}|$.
The orbit on the right of \eqref{flags} is always contained in the left,
and its stabilizer is the denominator subgroup by \eqref{ip}; consequently
this cardinality check really implies set equality.

We give the argument that this tests flag-transitivity, including flags
which might otherwise be omitted by a chamber-only model. Induct on the
size of a flag. Normalize its first $|J|$ objects to the base flag using
the induction hypothesis. Any further incident type-$k$ object lies in
$E(J,k)$ and hence can be normalized by the base-flag stabilizer
$\bigcap_{i\in J}P_i=H_{I\setminus J}$. Thus every flag is an image of a
base flag. In particular every flag extends to a chamber, and $H$ is
flag-transitive on the actual incidence system.

The chamber stabilizer is trivial by \eqref{ip}. A corank-one flag has
exactly two completions because its stabilizer is the two-element group
generated by its remaining mark. More generally, the stabilizer of a
base flag is generated by the remaining marks. Their chamber-adjacency
graph is connected. For a residue of rank at least two these successive
adjacent chambers share an object, so its incidence graph is connected.
The system is therefore a thin, residually connected geometry. Finally
any type-preserving automorphism fixing a chamber fixes its unique
adjacent chamber of each type and then every chamber. It is the identity.
Consequently the full type-preserving group is precisely $H$, and the
geometry is a regular hypertope.

All maximal subdiagrams here are finite. Their universal orders, in our
chosen vertex order, are respectively
\[
\begin{array}{ll}
3334:(24,48,48,24),&3335:(24,120,120,24),\\
3434:(48,48,48,48),&3534:(120,48,48,120),\\
353:(120,12,12,120),&435:(120,20,16,48),\\
535:(120,20,20,120),&5333:(8,24,120,120).
\end{array}
\]
These are the orders of $A_3,B_3,H_3$ or products of rank-one and
dihedral Coxeter groups. Equality with these orders makes the natural
surjection of each finite local Coxeter group injective. Otherwise the
actual local orders are retained explicitly in the witness records;
the broad column does not assume an injective local quotient.

\section{Exhaustion, source corrections, and retained premises}
Every finite regular hypertope with the prescribed diagram supplies its
distinguished involutions, hence an epimorphism $W(M)\to H$. Its kernel is
normal and has index $|H|$. This implication needs no injection on local
parabolics. Relabelling a source-equivalent diagram gives a marking of
our fixed presentation, so does not evade the normal-kernel search.
We enumerate all normal subgroups through the displayed bound, discard
quotients with wrong mark or pair orders, and apply \eqref{ip} and
\eqref{flags}. The only index-only prefilter is divisibility by
$\mathop{\rm lcm}_{i<j}(2M_{ij})$: each genuine rank-two marked subgroup
has order $2M_{ij}$. Thus no required quotient is lost.

The enumeration uses GAP 4.15.1 and LINS 0.9,
\texttt{LowIndexNormalSubgroupsSearchForAll}, with the unmodified default
options. Its documented contract is all normal subgroups through the
index bound, not one quotient per abstract isomorphism type. The installed
limit is $10^7$, checked before each search. The default has no early
termination or branch cut, and uses \texttt{UseLIS=false}. The algorithm
and its exact computer implementation are retained computational premises.
In particular LINS's tables of characteristically simple groups, Schur
multiplier primes, and transitive socle-overgroups are classification
inputs; we do not claim to reprove or regenerate their completeness.
The local documentation attributes the algorithm to Firth's 2005 thesis
and explains its pregenerated tables. Exact executable/table bytes and
clean terminal receipts are bound in the packet. This is not a formal
verification of GAP or LINS.

There are exactly three kernels of index336 for the first disputed square
and one of index2160 for the second. Every resulting marking fails
\eqref{ip}, even allowing smaller finite-diagram local residues.
Their compact tuples are retained and independently checked by literal
permutation closure. The source order $336$ and2160 rows therefore do not
describe regular hypertopes of their printed diagrams. A further negative
control is the order $360$ quotient of $(3,5,3,4)$: it satisfies all 256
intersections but has four, rather than two, type1 completions of a
particular rank-three base flag. It illustrates why local orders or the
intersection property alone cannot replace the incidence test.

There is also an elementary, more restricted obstruction at336: with
four universal $B_3$ residues, choose two maximal parabolics whose
remaining edge has label3. Their intersection would have order6, and
their product set would have $48^2/6=384$ elements. This argument alone
does \emph{not} exclude projective local residues; the unqualified
index336 exclusion above uses all three kernels.

Fresh and retained positive inventories are compared by a bijection of
ordered markings: an explicitly checked isomorphism sends each of the four
old marks to the corresponding new mark. This is equivalent to equality
of their kernels in the fixed presentation. Every record is matched once,
so multiplicities are preserved; equality of compact permutation bytes
is neither required nor presumed. The literal JSON inputs are also compared
with the exact GAP records before their independent audit.

The positive checks are repeated using independent standard-library
permutation closure and explicit coset sets. Neither compacting the GAP
actions nor printing structural descriptions supplies a completeness
premise. This component proves the finite assertions displayed above;
it does not finish Table 3's exact named extensions, the ninth square,
rank-five existence, the source's other classifications, or the whole
original. Initial universal-only checks and all capped square 3535
attempts are preserved as history and are not substituted for the broad
complete executions.

\paragraph{Attribution.}
The target rows, diagrams and minimum claims belong to
Fernandes--Leemans--Weiss, \emph{An Exploration of Locally Spherical
Regular Hypertopes}, fixed 16-page source, Sections 2 and 6, Tables 2--3.
The parent supplied the first local-order square 3434 probe and identified
the source's local-scope ambiguity. The present author supplied the full
incidence/intersection implementation, independent literal audit and
bounded core synthesis. A separate reader inspected the LINS documentation
and table boundary. None of these activities asserts priority for the
underlying finite-group algorithms or adds a new measurement or question.
\end{document}
